%% TREŚĆ ZADANIA, NUMER, IMIĘ I NAZWISKO

\begin{problem}{9.5}{Franciszek Pietrusiak}
Podaj przykład bijekcji:
\begin{multicols}{3}
\begin{enumerate}
    \item $\Z_+\to\Q_+$;
    \item $(0,1)\to\R$;
    \item $(0,1)\to[0,1]$.
\end{enumerate}
\end{multicols}
\end{problem}

%% MIEJSCE NA ZAWARTOŚĆ OPRACOWANIA
%% w razie potrzeby proszę zajrzeć do pliku 2-23.tex

\paragraph{Rozwiązanie} % można wpisać więcej niż jedno!

\subparagraph{(a)}
Rozpatrzmy drzewo binarne, w którym korzeń odpowiada ułamkowi $\frac{1}{1}$,
a dla dowolnego wierzchołka odpowiadającego dodatniemu ułamkowi $\frac{a}{b}$
zachodzi:
\begin{itemize}
    \item jego lewy syn odpowiada ułamkowi $\frac{a}{a+b}$
    \item jego prawy syn odpowiada ułamkowi $\frac{a+b}{b}$
\end{itemize}

Rysunek pierwszych trzech poziomów tak zdefiniowanego drzewa:

\begin{center}
\begin{tikzpicture}[
  level 1/.style={sibling distance=4cm, level distance=1cm},
  level 2/.style={sibling distance=2cm, level distance=1cm},
  every node/.style={align=center, inner sep=2pt, font=\small},
  edge from parent/.style={draw, ->, >=latex, thick}
]

\node {$\left(\frac{1}{1}\right)$ \\ {\color{maincolor}\scriptsize 1}}
  child { node {$\left(\frac{1}{2}\right)$ \\ {\color{maincolor}\scriptsize 2}}
    child { node {$\left(\frac{1}{3}\right)$ \\ {\color{maincolor}\scriptsize 4}} }
    child { node {$\left(\frac{3}{2}\right)$ \\ {\color{maincolor}\scriptsize 5}} }
  }
  child { node {$\left(\frac{2}{1}\right)$ \\ {\color{maincolor}\scriptsize 3}}
    child { node {$\left(\frac{2}{3}\right)$ \\ {\color{maincolor}\scriptsize 6}} }
    child { node {$\left(\frac{3}{1}\right)$ \\ {\color{maincolor}\scriptsize 7}} }
  };
\end{tikzpicture}
\end{center}

Pod dowolnym wierzchołkiem $\left(\frac{p}{q}\right)$ kolorem
\textcolor{maincolor}{zielonym}
zaznaczono jego numer $n_{\left(\frac{p}{q}\right)} \in \Z_+$.
Numery wierzchołków rosną w pierwszej kolejności po poziomie, a
następnie "od lewej do prawej". Takie numerowanie jest oczywiście
jednoznaczne.

Wtedy szukana bijekcja jest przyporządkowaniem, które
dla każdej liczby $n \in Z_+$ zwraca ułamek $\frac{p}{q} \in \Q_+$
który odpowiada wierzchołkowi $\left(\frac{p}{q}\right)$ t.że
$n_{\left(\frac{p}{q}\right)} = n$.

Pozostało udowodnić, że:
\begin{itemize}
    \item ułamki w drzewie są nieskracalne
    \item nie ma powtórzeń (różnowartościowość)
    \item każdy ułamek znajduje się w drzewie ("na")
\end{itemize}

\begin{proof}[Dowód (pierwszy punkt)]
Baza Indukcji: Dla korzenia mamy $NWD(1, 1) = 1$.

Krok Indukcyjny: Załóżmy, że $\nwd(a, b) = 1$.
Wtedy na mocy \textit{Algorytmu Euklidesa}:
$\nwd(a, a+b) = \nwd(a, b) = 1$
oraz analogicznie
$\nwd(a+b, b) = \nwd(a, b) = 1$.
Czyli dzieci wierzchołka $\left(\frac{a}{b}\right)$
odpowiadają ułamkom nieskracalnym.

Z konstrukcji drzewa wynika więc,
że wszystkie wierzchołki odpowiadają ułamkom nieskracalnym.
\end{proof}

\begin{proof}[Dowód (drugi punkt)]
Przypomnijmy zasady tworzenia drzewa:
Dla dowolnego wierzchołka $\left(\frac{a}{b}\right)$ jego
lewy syn to $\left(\frac{a}{a+b}\right)$, a prawy $\left(\frac{a+b}{b}\right)$.
Skoro $a, b > 0$ to lewy syn odpowiada ułamkowi mniejszemu
od $1$. Prawy natomiast odpowiada ułamkowi większemu od $1$.

Zatem biorąc dowolny wierzchołek $\left(\frac{p}{q}\right)$
jednoznacznie wiadomo którym jest synem.
Jeśli lewym, to jego ojcem jest $\left(\frac{p}{q-p}\right)$,
a jeśli prawym to ojcem jest $\left(\frac{p-q}{q}\right)$.

To oznacza, że dla każdego wierzchołka istnieje jednoznacznie
wyznaczona ścieżka idąca od niego do korzenia drzewa. Co za tym
idzie nie istnieją dwa różne wierzchołki którym odpowiada ten
sam ułamek.
\end{proof}

\begin{proof}[Dowód (trzeci punkt)]
Weźmy dowolną dodatnią liczbę wymierną i zapiszmy ją jako nieskracalny
ułamek $\frac{p}{q}$. Odpowiadać mu będzie wierzchołek 
$v := \left(\frac{p}{q}\right)$.

Z punktu drugiego wiemy, że jeśli $p < q$ to $v$ jest lewym synem
i jego ojcem jest $\left(\frac{p}{q-p}\right)$, a jeśli $p > q$
to $v$ prawym synem, a jego ojcem jest $\left(\frac{p-q}{q}\right)$.

W obu przypadkach przechodząc z syna do ojca suma licznika i mianownika
się zmniejsza. Gdy $v$ jest lewym synem: $p+q \to p+(q-p)=q$, gdy prawym:
$p+q \to (p-q)+q=p$. Skoro jednak wszystkie wierzchołki w drzewie odpowiadają
ułamkom dodatnim to ta suma nie może maleć w nieskończoność. Proces
przechodzenia z syna do ojca zatrzyma się wtedy gdy $p = q$.
Skoro $\nwd(p, q) = 1$ to $\left(\frac{p}{q}\right) = \left(\frac{1}{1}\right)$
czyli proces zatrzyma się w korzeniu drzewa. Właśnie pokazaliśmy, że
istnieje ścieżka z $v$ do korzenia drzewa, a to oczywiście oznacza,
że $v$ należy do drzewa. Z dowolności wyboru ułamka wynika teza punktu trzeciego.
\end{proof}

\subparagraph{(b)}
Zacznijmy od funkcji $f_1: \left(-\frac{\pi}{2}, \frac{\pi}{2}\right) \to \R$
określonej wzorem $f_1(x) = \tg(x)$.

Z własności funkcji tangens wiemy, że na tym zbiorze
argumentów $f_1$ jest funkcją ściśle rosnącą $\implies$
$f_1$ jest funkcją różnowartościową.

Dodatkowo wiadomo, że $f_1$ jest funkcją ciągłą oraz
$$\lim_{x \to (-\frac{\pi}{2})^+} f_1(x) = -\infty \quad \lim_{x \to (\frac{\pi}{2})^-} f_1(x) = +\infty$$
więc $f_1$ przyjmuje wszystkie wartości z $\R$ więc jest funkcją "na".

Skoro $f_1$ jest zarówno funkcją różnowartościową jak i funkcją "na",
to jest ona bijekcją.

Jednak w naszym zadaniu potrzebujemy bijekcji, której dziedzina to $(0,1)$.

Aby to naprawić zdefiniujemy dwie dodatkowe funkcje:

Niech $f_2: \left(-\frac{1}{2}, \frac{1}{2}\right) \to \left(-\frac{\pi}{2}, \frac{\pi}{2}\right)$
t.że $f_2(x) = \pi x$
oraz
$f_3: (0,1) \to \left(-\frac{1}{2}, \frac{1}{2}\right)$ t.że $f_3(x) = x - \frac{1}{2}$.
Nietrudno zauważyć, że te funkcje są zarówno różnowartościowe jak i "na"
$\implies$ są to bijekcje.

Wtedy funkcja
$$f(x) = (f_1 \circ f_2 \circ f_3)(x) = f_1(f_2(f_3(x))) = \tg\left(\pi \left(x - \frac{1}{2}\right)\right)$$
będąca złożeniem trzech bijekcji jest szukaną bijekcją.


% \paragraph{Wskazówka} % może być więcej niż jedna!
\subparagraph{Wskazówka 1. (b)}
Pomyśl o funkcjach, które mają asymptoty pionowe.

\paragraph{Komentarz} % pogłębiony komentarz metodyczny dotyczący zadania

\paragraph{Uogólnienia i modyfikacje} % jeśli są naturalne, np. łatwiejsze / trudniejsze wersje zadania nadające się do pracy na różnych poziomach edukacyjnych

% można dopisywać kolejne akapity, np. ciekawostki, możliwe dygresje, linki do wartościowych źródeł itp.